% ==============================================================
%  radar_styles.tex
%  Core palette, libraries, and primitives for the Radar Night
%  TikZ graphics library (radar processing / geometry / NN icons)
% ==============================================================
%
% Usage:
%   % ==============================================================
%  radar_styles.tex
%  Core palette, libraries, and primitives for the Radar Night
%  TikZ graphics library (radar processing / geometry / NN icons)
% ==============================================================
%
% Usage:
%   % ==============================================================
%  radar_styles.tex
%  Core palette, libraries, and primitives for the Radar Night
%  TikZ graphics library (radar processing / geometry / NN icons)
% ==============================================================
%
% Usage:
%   % ==============================================================
%  radar_styles.tex
%  Core palette, libraries, and primitives for the Radar Night
%  TikZ graphics library (radar processing / geometry / NN icons)
% ==============================================================
%
% Usage:
%   \input{radar_styles.tex}
% or, if used as a package-like include inside a larger project,
% simply \input this file once before any of the icon files.

% ---------------- required TikZ libraries ---------------------
\usetikzlibrary{
  shapes.multipart,
  shapes.geometric,
  shapes.arrows,
  arrows.meta,
  positioning,
  calc,
  fit,
  decorations.pathmorphing,
  decorations.pathreplacing,
  patterns,
  patterns.meta,
  shadows,
  backgrounds,
  3d
}

% ------------------------------------------------------------
% 1. COLOR PALETTE  ("Radar Night" -- matches all diagram assets)
% ------------------------------------------------------------
\definecolor{radarNavy}{HTML}{1A2332}
\definecolor{radarPurple}{HTML}{6B4E9E}
\definecolor{radarOrange}{HTML}{E8734A}
\definecolor{radarBlue}{HTML}{3B7A9E}
\definecolor{radarGreen}{HTML}{4A9B6E}
\definecolor{radarBG}{HTML}{FAFAF8}
\definecolor{radarGrey}{HTML}{5A6472}
\definecolor{kuldefault}{HTML}{00407a}
\definecolor{kuldark}{HTML}{123C75}    % heading background
\definecolor{kullight}{HTML}{F0F7FA}   % block background
\definecolor{kulbg}{RGB}{29,141,176}   % background fill
\definecolor{kulblue}{HTML}{00407a}

\colorlet{blocktitlebgcolor}{radarBlue}
\colorlet{blocktitlefgcolor}{white}
\colorlet{blockbodybgcolor}{white}
\colorlet{backgroundcolor}{radarBG}
\colorlet{framecolor}{radarNavy}

% Semantic aliases used across the library so individual icon
% files never hard-code a raw color -- only these roles.
\colorlet{radarProcessing}{radarBlue}   % radar-domain blocks
\colorlet{deepLearning}{radarGreen}     % NN / learning blocks
\colorlet{resultsColor}{radarOrange}    % results / metrics
\colorlet{geometryColor}{radarPurple}   % geometry / point-cloud blocks
\colorlet{iconStroke}{radarNavy}
\colorlet{iconFaint}{radarGrey!35}
\colorlet{cardBorder}{radarNavy}
\colorlet{cardBG}{white}

% ------------------------------------------------------------
% 2. GLOBAL PGFKEYS  (per-module color role, size, etc.)
% ------------------------------------------------------------
% Every icon/card macro in this library accepts a `color=` key
% that maps to one of the semantic roles above, so a whole
% figure can be recolored by module without touching the icons.
\pgfkeys{
  /radarlib/.is family, /radarlib,
  default/.style = {
    color       = radarProcessing,
    scale       = 1,
    width       = 2.1cm,
  },
  color/.estore in      = \radarlibColor,
  scale/.estore in      = \radarlibScale,
  width/.estore in      = \radarlibWidth,
}
\pgfkeys{/radarlib, default}

% ------------------------------------------------------------
% 3. FONTS / TEXT STYLES
% ------------------------------------------------------------
\tikzset{
  card title/.style   = {font=\sffamily\bfseries\footnotesize, text=blocktitlefgcolor},
  card subtitle/.style = {font=\sffamily\footnotesize, text=radarGrey},
  metric number/.style = {font=\sffamily\bfseries\Large, text=radarNavy},
  metric label/.style  = {font=\sffamily\footnotesize, text=radarGrey},
}
% Plain-text equivalents of the styles above, for use *inside*
% already-open node text (where a tikz style key can't be
% invoked directly) -- e.g. {\metricNumberFont 73.8\%}.
\newcommand{\cardTitleFont}{\sffamily\bfseries\footnotesize\color{blocktitlefgcolor}}
\newcommand{\cardSubtitleFont}{\sffamily\footnotesize\color{radarGrey}}
\newcommand{\metricNumberFont}{\sffamily\bfseries\Large\color{radarNavy}}
\newcommand{\metricLabelFont}{\sffamily\footnotesize\color{radarGrey}}

% ------------------------------------------------------------
% 4. FLOW ARROW  (the connector used between every pipeline stage)
% ------------------------------------------------------------
% Usage:  \draw[flow] (a) -- (b);
%         \draw[flow, color=deepLearning] (a) -- (b);
\tikzset{
  flow/.style = {
    -{Latex[length=2.6mm, width=2mm]},
    line width=0.9pt,
    draw=radarNavy,
    shorten <=1pt,
    shorten >=1pt,
  },
  flow dashed/.style = {
    flow,
    dash pattern=on 2.5pt off 2pt,
  },
}

% Convenience macro so a simple horizontal/vertical chain reads
% declaratively: \PipelineArrow{a}{b}
\newcommand{\PipelineArrow}[3][flow]{\draw[#1] (#2) -- (#3);}

% ------------------------------------------------------------
% 5. PROCESS CARD
% ------------------------------------------------------------
% The workhorse of the whole poster: an icon on top, a colored
% title bar on the bottom, rounded corners, subtle border.
% Built on the `rectangle split` shape so the icon and title
% live in one node (keeps alignment automatic).
%
% Usage inside a tikzpicture:
%   \ProcessCard{a}{Raw Radar}{\RadarCube}
%   \ProcessCard[right=1.4cm of a, titlecolor=deepLearning]{b}{Range FFT}{\RangeFFT}
%
% Args: #1 = extra node options (positioning etc, optional)
%       #2 = node name
%       #3 = title text
%       #4 = icon content (TikZ draw commands, drawn in a small
%            nested picture with origin at the card center)
\tikzset{
  process card/.style = {
    rectangle split,
    rectangle split parts=2,
    rectangle split part align={center, center},
    draw=cardBorder,
    line width=0.8pt,
    rounded corners=4pt,
    text width=2.1cm,
    align=center,
    inner sep=5pt,
    drop shadow={shadow xshift=0.6pt, shadow yshift=-0.6pt, opacity=0.25},
    rectangle split part fill={cardBG, radarProcessing},
  },
  % Safe way to recolor a single card's title bar, e.g.:
  %   \ProcessCard[card color=deepLearning]{b}{Range FFT}{...}
  % (Do not set `fill=` directly on a process card -- it conflicts
  % with `rectangle split part fill` and blanks the node.)
  card color/.style = {rectangle split part fill={cardBG, #1}},
}

\newcommand{\ProcessCard}[4][]{%
  \node[process card, #1] (#2) {%
    \nodepart{one}
      \begin{tikzpicture}[baseline, remember picture, every node/.style={shape=rectangle}]
        #4
      \end{tikzpicture}
    \nodepart{two}
      \textcolor{blocktitlefgcolor}{\sffamily\bfseries\footnotesize #3}
  };%
}

% ------------------------------------------------------------
% 6. METRIC BOX  (the big "73.8% IoU" style callouts)
% ------------------------------------------------------------
% Usage: \MetricBox{a}{0,0}{73.8\%}{IoU}
\newcommand{\MetricBox}[4]{%
  \node[align=center, inner sep=2pt] (#1) at (#2) {%
    {\metricNumberFont #3}\\[1pt]
    {\metricLabelFont #4}%
  };%
}

% ------------------------------------------------------------
% 7. LEGEND BOX  (small colored swatch + label, for figure legends)
% ------------------------------------------------------------
% Usage: \LegendSwatch{0,0}{radarBlue}{Radar processing}
\newcommand{\LegendSwatch}[3]{%
  \begin{scope}[shift={(#1)}]
    \node[fill=#2, rounded corners=1.5pt, minimum width=10pt, minimum height=10pt, anchor=west] (sw) {};
    \node[anchor=west, font=\sffamily\footnotesize, text=radarGrey, right=4pt of sw] {#3};
  \end{scope}
}

% ------------------------------------------------------------
% 8. SECTION HEADER  (poster column headers, e.g. "Contribution 1")
% ------------------------------------------------------------
\newcommand{\SectionHeader}[2][radarProcessing]{%
  \node[fill=#1, text=white, rounded corners=2pt, inner sep=6pt,
        font=\sffamily\bfseries\normalsize, minimum width=4cm] {#2};%
}

% ------------------------------------------------------------
% 9. RESULTS TABLE  (palette-matched booktabs-style table)
% ------------------------------------------------------------
% Not a TikZ node -- a small LaTeX environment wrapper so results
% tables share the same fonts/colors/weight as the figures. Needs
% \usepackage{booktabs} and \usepackage[table]{xcolor} in the
% poster/paper preamble.
%
% Usage:
%   \begin{ResultsTable}{lcc}
%     Method          & Supervised & Few-shot \\
%     TCPCN           & 96.2       & 35.1     \\
%     PCTtrans        & 99.0       & 32.0     \\
%     \BestRow Ours + Pretrain & 94.4 & 70.0  \\
%   \end{ResultsTable}
\newenvironment{ResultsTable}[1]{%
  \renewcommand{\arraystretch}{1.35}%
  \begin{tabular}{#1}
  \toprule
}{%
  \bottomrule
  \end{tabular}%
}
% Header row: white bold text on the radarProcessing band. Formatting
% resets after \midrule so it doesn't bleed into the data rows below.
\newcommand{\HeaderRow}[1]{%
  \rowcolor{radarProcessing}\bfseries\color{white}#1\\
  \midrule
  \normalfont\color{black}%
}
% Highlight a single data row (e.g. your best result) in results color.
\newcommand{\BestRow}{\rowcolor{resultsColor!25}}

% ------------------------------------------------------------
% 10. LEGEND BOX (complete)
% ------------------------------------------------------------
% Usage: \LegendBox{x}{y}{color}{label}
\newcommand{\LegendBox}[4]{%
  \begin{scope}[shift={(#1,#2)}]
    \node[fill=#3, rounded corners=1.5pt, minimum width=8pt, minimum height=8pt, anchor=west] (sw) at (0,0) {};
    \node[anchor=west, font=\sffamily\footnotesize, text=radarGrey, right=6pt of sw] {#4};
  \end{scope}
}
% or, if used as a package-like include inside a larger project,
% simply \input this file once before any of the icon files.

% ---------------- required TikZ libraries ---------------------
\usetikzlibrary{
  shapes.multipart,
  shapes.geometric,
  shapes.arrows,
  arrows.meta,
  positioning,
  calc,
  fit,
  decorations.pathmorphing,
  decorations.pathreplacing,
  patterns,
  patterns.meta,
  shadows,
  backgrounds,
  3d
}

% ------------------------------------------------------------
% 1. COLOR PALETTE  ("Radar Night" -- matches all diagram assets)
% ------------------------------------------------------------
\definecolor{radarNavy}{HTML}{1A2332}
\definecolor{radarPurple}{HTML}{6B4E9E}
\definecolor{radarOrange}{HTML}{E8734A}
\definecolor{radarBlue}{HTML}{3B7A9E}
\definecolor{radarGreen}{HTML}{4A9B6E}
\definecolor{radarBG}{HTML}{FAFAF8}
\definecolor{radarGrey}{HTML}{5A6472}
\definecolor{kuldefault}{HTML}{00407a}
\definecolor{kuldark}{HTML}{123C75}    % heading background
\definecolor{kullight}{HTML}{F0F7FA}   % block background
\definecolor{kulbg}{RGB}{29,141,176}   % background fill
\definecolor{kulblue}{HTML}{00407a}

\colorlet{blocktitlebgcolor}{radarBlue}
\colorlet{blocktitlefgcolor}{white}
\colorlet{blockbodybgcolor}{white}
\colorlet{backgroundcolor}{radarBG}
\colorlet{framecolor}{radarNavy}

% Semantic aliases used across the library so individual icon
% files never hard-code a raw color -- only these roles.
\colorlet{radarProcessing}{radarBlue}   % radar-domain blocks
\colorlet{deepLearning}{radarGreen}     % NN / learning blocks
\colorlet{resultsColor}{radarOrange}    % results / metrics
\colorlet{geometryColor}{radarPurple}   % geometry / point-cloud blocks
\colorlet{iconStroke}{radarNavy}
\colorlet{iconFaint}{radarGrey!35}
\colorlet{cardBorder}{radarNavy}
\colorlet{cardBG}{white}

% ------------------------------------------------------------
% 2. GLOBAL PGFKEYS  (per-module color role, size, etc.)
% ------------------------------------------------------------
% Every icon/card macro in this library accepts a `color=` key
% that maps to one of the semantic roles above, so a whole
% figure can be recolored by module without touching the icons.
\pgfkeys{
  /radarlib/.is family, /radarlib,
  default/.style = {
    color       = radarProcessing,
    scale       = 1,
    width       = 2.1cm,
  },
  color/.estore in      = \radarlibColor,
  scale/.estore in      = \radarlibScale,
  width/.estore in      = \radarlibWidth,
}
\pgfkeys{/radarlib, default}

% ------------------------------------------------------------
% 3. FONTS / TEXT STYLES
% ------------------------------------------------------------
\tikzset{
  card title/.style   = {font=\sffamily\bfseries\footnotesize, text=blocktitlefgcolor},
  card subtitle/.style = {font=\sffamily\footnotesize, text=radarGrey},
  metric number/.style = {font=\sffamily\bfseries\Large, text=radarNavy},
  metric label/.style  = {font=\sffamily\footnotesize, text=radarGrey},
}
% Plain-text equivalents of the styles above, for use *inside*
% already-open node text (where a tikz style key can't be
% invoked directly) -- e.g. {\metricNumberFont 73.8\%}.
\newcommand{\cardTitleFont}{\sffamily\bfseries\footnotesize\color{blocktitlefgcolor}}
\newcommand{\cardSubtitleFont}{\sffamily\footnotesize\color{radarGrey}}
\newcommand{\metricNumberFont}{\sffamily\bfseries\Large\color{radarNavy}}
\newcommand{\metricLabelFont}{\sffamily\footnotesize\color{radarGrey}}

% ------------------------------------------------------------
% 4. FLOW ARROW  (the connector used between every pipeline stage)
% ------------------------------------------------------------
% Usage:  \draw[flow] (a) -- (b);
%         \draw[flow, color=deepLearning] (a) -- (b);
\tikzset{
  flow/.style = {
    -{Latex[length=2.6mm, width=2mm]},
    line width=0.9pt,
    draw=radarNavy,
    shorten <=1pt,
    shorten >=1pt,
  },
  flow dashed/.style = {
    flow,
    dash pattern=on 2.5pt off 2pt,
  },
}

% Convenience macro so a simple horizontal/vertical chain reads
% declaratively: \PipelineArrow{a}{b}
\newcommand{\PipelineArrow}[3][flow]{\draw[#1] (#2) -- (#3);}

% ------------------------------------------------------------
% 5. PROCESS CARD
% ------------------------------------------------------------
% The workhorse of the whole poster: an icon on top, a colored
% title bar on the bottom, rounded corners, subtle border.
% Built on the `rectangle split` shape so the icon and title
% live in one node (keeps alignment automatic).
%
% Usage inside a tikzpicture:
%   \ProcessCard{a}{Raw Radar}{\RadarCube}
%   \ProcessCard[right=1.4cm of a, titlecolor=deepLearning]{b}{Range FFT}{\RangeFFT}
%
% Args: #1 = extra node options (positioning etc, optional)
%       #2 = node name
%       #3 = title text
%       #4 = icon content (TikZ draw commands, drawn in a small
%            nested picture with origin at the card center)
\tikzset{
  process card/.style = {
    rectangle split,
    rectangle split parts=2,
    rectangle split part align={center, center},
    draw=cardBorder,
    line width=0.8pt,
    rounded corners=4pt,
    text width=2.1cm,
    align=center,
    inner sep=5pt,
    drop shadow={shadow xshift=0.6pt, shadow yshift=-0.6pt, opacity=0.25},
    rectangle split part fill={cardBG, radarProcessing},
  },
  % Safe way to recolor a single card's title bar, e.g.:
  %   \ProcessCard[card color=deepLearning]{b}{Range FFT}{...}
  % (Do not set `fill=` directly on a process card -- it conflicts
  % with `rectangle split part fill` and blanks the node.)
  card color/.style = {rectangle split part fill={cardBG, #1}},
}

\newcommand{\ProcessCard}[4][]{%
  \node[process card, #1] (#2) {%
    \nodepart{one}
      \begin{tikzpicture}[baseline, remember picture, every node/.style={shape=rectangle}]
        #4
      \end{tikzpicture}
    \nodepart{two}
      \textcolor{blocktitlefgcolor}{\sffamily\bfseries\footnotesize #3}
  };%
}

% ------------------------------------------------------------
% 6. METRIC BOX  (the big "73.8% IoU" style callouts)
% ------------------------------------------------------------
% Usage: \MetricBox{a}{0,0}{73.8\%}{IoU}
\newcommand{\MetricBox}[4]{%
  \node[align=center, inner sep=2pt] (#1) at (#2) {%
    {\metricNumberFont #3}\\[1pt]
    {\metricLabelFont #4}%
  };%
}

% ------------------------------------------------------------
% 7. LEGEND BOX  (small colored swatch + label, for figure legends)
% ------------------------------------------------------------
% Usage: \LegendSwatch{0,0}{radarBlue}{Radar processing}
\newcommand{\LegendSwatch}[3]{%
  \begin{scope}[shift={(#1)}]
    \node[fill=#2, rounded corners=1.5pt, minimum width=10pt, minimum height=10pt, anchor=west] (sw) {};
    \node[anchor=west, font=\sffamily\footnotesize, text=radarGrey, right=4pt of sw] {#3};
  \end{scope}
}

% ------------------------------------------------------------
% 8. SECTION HEADER  (poster column headers, e.g. "Contribution 1")
% ------------------------------------------------------------
\newcommand{\SectionHeader}[2][radarProcessing]{%
  \node[fill=#1, text=white, rounded corners=2pt, inner sep=6pt,
        font=\sffamily\bfseries\normalsize, minimum width=4cm] {#2};%
}

% ------------------------------------------------------------
% 9. RESULTS TABLE  (palette-matched booktabs-style table)
% ------------------------------------------------------------
% Not a TikZ node -- a small LaTeX environment wrapper so results
% tables share the same fonts/colors/weight as the figures. Needs
% \usepackage{booktabs} and \usepackage[table]{xcolor} in the
% poster/paper preamble.
%
% Usage:
%   \begin{ResultsTable}{lcc}
%     Method          & Supervised & Few-shot \\
%     TCPCN           & 96.2       & 35.1     \\
%     PCTtrans        & 99.0       & 32.0     \\
%     \BestRow Ours + Pretrain & 94.4 & 70.0  \\
%   \end{ResultsTable}
\newenvironment{ResultsTable}[1]{%
  \renewcommand{\arraystretch}{1.35}%
  \begin{tabular}{#1}
  \toprule
}{%
  \bottomrule
  \end{tabular}%
}
% Header row: white bold text on the radarProcessing band. Formatting
% resets after \midrule so it doesn't bleed into the data rows below.
\newcommand{\HeaderRow}[1]{%
  \rowcolor{radarProcessing}\bfseries\color{white}#1\\
  \midrule
  \normalfont\color{black}%
}
% Highlight a single data row (e.g. your best result) in results color.
\newcommand{\BestRow}{\rowcolor{resultsColor!25}}

% ------------------------------------------------------------
% 10. LEGEND BOX (complete)
% ------------------------------------------------------------
% Usage: \LegendBox{x}{y}{color}{label}
\newcommand{\LegendBox}[4]{%
  \begin{scope}[shift={(#1,#2)}]
    \node[fill=#3, rounded corners=1.5pt, minimum width=8pt, minimum height=8pt, anchor=west] (sw) at (0,0) {};
    \node[anchor=west, font=\sffamily\footnotesize, text=radarGrey, right=6pt of sw] {#4};
  \end{scope}
}
% or, if used as a package-like include inside a larger project,
% simply \input this file once before any of the icon files.

% ---------------- required TikZ libraries ---------------------
\usetikzlibrary{
  shapes.multipart,
  shapes.geometric,
  shapes.arrows,
  arrows.meta,
  positioning,
  calc,
  fit,
  decorations.pathmorphing,
  decorations.pathreplacing,
  patterns,
  patterns.meta,
  shadows,
  backgrounds,
  3d
}

% ------------------------------------------------------------
% 1. COLOR PALETTE  ("Radar Night" -- matches all diagram assets)
% ------------------------------------------------------------
\definecolor{radarNavy}{HTML}{1A2332}
\definecolor{radarPurple}{HTML}{6B4E9E}
\definecolor{radarOrange}{HTML}{E8734A}
\definecolor{radarBlue}{HTML}{3B7A9E}
\definecolor{radarGreen}{HTML}{4A9B6E}
\definecolor{radarBG}{HTML}{FAFAF8}
\definecolor{radarGrey}{HTML}{5A6472}
\definecolor{kuldefault}{HTML}{00407a}
\definecolor{kuldark}{HTML}{123C75}    % heading background
\definecolor{kullight}{HTML}{F0F7FA}   % block background
\definecolor{kulbg}{RGB}{29,141,176}   % background fill
\definecolor{kulblue}{HTML}{00407a}

\colorlet{blocktitlebgcolor}{radarBlue}
\colorlet{blocktitlefgcolor}{white}
\colorlet{blockbodybgcolor}{white}
\colorlet{backgroundcolor}{radarBG}
\colorlet{framecolor}{radarNavy}

% Semantic aliases used across the library so individual icon
% files never hard-code a raw color -- only these roles.
\colorlet{radarProcessing}{radarBlue}   % radar-domain blocks
\colorlet{deepLearning}{radarGreen}     % NN / learning blocks
\colorlet{resultsColor}{radarOrange}    % results / metrics
\colorlet{geometryColor}{radarPurple}   % geometry / point-cloud blocks
\colorlet{iconStroke}{radarNavy}
\colorlet{iconFaint}{radarGrey!35}
\colorlet{cardBorder}{radarNavy}
\colorlet{cardBG}{white}

% ------------------------------------------------------------
% 2. GLOBAL PGFKEYS  (per-module color role, size, etc.)
% ------------------------------------------------------------
% Every icon/card macro in this library accepts a `color=` key
% that maps to one of the semantic roles above, so a whole
% figure can be recolored by module without touching the icons.
\pgfkeys{
  /radarlib/.is family, /radarlib,
  default/.style = {
    color       = radarProcessing,
    scale       = 1,
    width       = 2.1cm,
  },
  color/.estore in      = \radarlibColor,
  scale/.estore in      = \radarlibScale,
  width/.estore in      = \radarlibWidth,
}
\pgfkeys{/radarlib, default}

% ------------------------------------------------------------
% 3. FONTS / TEXT STYLES
% ------------------------------------------------------------
\tikzset{
  card title/.style   = {font=\sffamily\bfseries\footnotesize, text=blocktitlefgcolor},
  card subtitle/.style = {font=\sffamily\footnotesize, text=radarGrey},
  metric number/.style = {font=\sffamily\bfseries\Large, text=radarNavy},
  metric label/.style  = {font=\sffamily\footnotesize, text=radarGrey},
}
% Plain-text equivalents of the styles above, for use *inside*
% already-open node text (where a tikz style key can't be
% invoked directly) -- e.g. {\metricNumberFont 73.8\%}.
\newcommand{\cardTitleFont}{\sffamily\bfseries\footnotesize\color{blocktitlefgcolor}}
\newcommand{\cardSubtitleFont}{\sffamily\footnotesize\color{radarGrey}}
\newcommand{\metricNumberFont}{\sffamily\bfseries\Large\color{radarNavy}}
\newcommand{\metricLabelFont}{\sffamily\footnotesize\color{radarGrey}}

% ------------------------------------------------------------
% 4. FLOW ARROW  (the connector used between every pipeline stage)
% ------------------------------------------------------------
% Usage:  \draw[flow] (a) -- (b);
%         \draw[flow, color=deepLearning] (a) -- (b);
\tikzset{
  flow/.style = {
    -{Latex[length=2.6mm, width=2mm]},
    line width=0.9pt,
    draw=radarNavy,
    shorten <=1pt,
    shorten >=1pt,
  },
  flow dashed/.style = {
    flow,
    dash pattern=on 2.5pt off 2pt,
  },
}

% Convenience macro so a simple horizontal/vertical chain reads
% declaratively: \PipelineArrow{a}{b}
\newcommand{\PipelineArrow}[3][flow]{\draw[#1] (#2) -- (#3);}

% ------------------------------------------------------------
% 5. PROCESS CARD
% ------------------------------------------------------------
% The workhorse of the whole poster: an icon on top, a colored
% title bar on the bottom, rounded corners, subtle border.
% Built on the `rectangle split` shape so the icon and title
% live in one node (keeps alignment automatic).
%
% Usage inside a tikzpicture:
%   \ProcessCard{a}{Raw Radar}{\RadarCube}
%   \ProcessCard[right=1.4cm of a, titlecolor=deepLearning]{b}{Range FFT}{\RangeFFT}
%
% Args: #1 = extra node options (positioning etc, optional)
%       #2 = node name
%       #3 = title text
%       #4 = icon content (TikZ draw commands, drawn in a small
%            nested picture with origin at the card center)
\tikzset{
  process card/.style = {
    rectangle split,
    rectangle split parts=2,
    rectangle split part align={center, center},
    draw=cardBorder,
    line width=0.8pt,
    rounded corners=4pt,
    text width=2.1cm,
    align=center,
    inner sep=5pt,
    drop shadow={shadow xshift=0.6pt, shadow yshift=-0.6pt, opacity=0.25},
    rectangle split part fill={cardBG, radarProcessing},
  },
  % Safe way to recolor a single card's title bar, e.g.:
  %   \ProcessCard[card color=deepLearning]{b}{Range FFT}{...}
  % (Do not set `fill=` directly on a process card -- it conflicts
  % with `rectangle split part fill` and blanks the node.)
  card color/.style = {rectangle split part fill={cardBG, #1}},
}

\newcommand{\ProcessCard}[4][]{%
  \node[process card, #1] (#2) {%
    \nodepart{one}
      \begin{tikzpicture}[baseline, remember picture, every node/.style={shape=rectangle}]
        #4
      \end{tikzpicture}
    \nodepart{two}
      \textcolor{blocktitlefgcolor}{\sffamily\bfseries\footnotesize #3}
  };%
}

% ------------------------------------------------------------
% 6. METRIC BOX  (the big "73.8% IoU" style callouts)
% ------------------------------------------------------------
% Usage: \MetricBox{a}{0,0}{73.8\%}{IoU}
\newcommand{\MetricBox}[4]{%
  \node[align=center, inner sep=2pt] (#1) at (#2) {%
    {\metricNumberFont #3}\\[1pt]
    {\metricLabelFont #4}%
  };%
}

% ------------------------------------------------------------
% 7. LEGEND BOX  (small colored swatch + label, for figure legends)
% ------------------------------------------------------------
% Usage: \LegendSwatch{0,0}{radarBlue}{Radar processing}
\newcommand{\LegendSwatch}[3]{%
  \begin{scope}[shift={(#1)}]
    \node[fill=#2, rounded corners=1.5pt, minimum width=10pt, minimum height=10pt, anchor=west] (sw) {};
    \node[anchor=west, font=\sffamily\footnotesize, text=radarGrey, right=4pt of sw] {#3};
  \end{scope}
}

% ------------------------------------------------------------
% 8. SECTION HEADER  (poster column headers, e.g. "Contribution 1")
% ------------------------------------------------------------
\newcommand{\SectionHeader}[2][radarProcessing]{%
  \node[fill=#1, text=white, rounded corners=2pt, inner sep=6pt,
        font=\sffamily\bfseries\normalsize, minimum width=4cm] {#2};%
}

% ------------------------------------------------------------
% 9. RESULTS TABLE  (palette-matched booktabs-style table)
% ------------------------------------------------------------
% Not a TikZ node -- a small LaTeX environment wrapper so results
% tables share the same fonts/colors/weight as the figures. Needs
% \usepackage{booktabs} and \usepackage[table]{xcolor} in the
% poster/paper preamble.
%
% Usage:
%   \begin{ResultsTable}{lcc}
%     Method          & Supervised & Few-shot \\
%     TCPCN           & 96.2       & 35.1     \\
%     PCTtrans        & 99.0       & 32.0     \\
%     \BestRow Ours + Pretrain & 94.4 & 70.0  \\
%   \end{ResultsTable}
\newenvironment{ResultsTable}[1]{%
  \renewcommand{\arraystretch}{1.35}%
  \begin{tabular}{#1}
  \toprule
}{%
  \bottomrule
  \end{tabular}%
}
% Header row: white bold text on the radarProcessing band. Formatting
% resets after \midrule so it doesn't bleed into the data rows below.
\newcommand{\HeaderRow}[1]{%
  \rowcolor{radarProcessing}\bfseries\color{white}#1\\
  \midrule
  \normalfont\color{black}%
}
% Highlight a single data row (e.g. your best result) in results color.
\newcommand{\BestRow}{\rowcolor{resultsColor!25}}

% ------------------------------------------------------------
% 10. LEGEND BOX (complete)
% ------------------------------------------------------------
% Usage: \LegendBox{x}{y}{color}{label}
\newcommand{\LegendBox}[4]{%
  \begin{scope}[shift={(#1,#2)}]
    \node[fill=#3, rounded corners=1.5pt, minimum width=8pt, minimum height=8pt, anchor=west] (sw) at (0,0) {};
    \node[anchor=west, font=\sffamily\footnotesize, text=radarGrey, right=6pt of sw] {#4};
  \end{scope}
}
% or, if used as a package-like include inside a larger project,
% simply \input this file once before any of the icon files.

% ---------------- required TikZ libraries ---------------------
\usetikzlibrary{
  shapes.multipart,
  shapes.geometric,
  shapes.arrows,
  arrows.meta,
  positioning,
  calc,
  fit,
  decorations.pathmorphing,
  decorations.pathreplacing,
  patterns,
  patterns.meta,
  shadows,
  backgrounds,
  3d
}

% ------------------------------------------------------------
% 1. COLOR PALETTE  ("Radar Night" -- matches all diagram assets)
% ------------------------------------------------------------
\definecolor{radarNavy}{HTML}{1A2332}
\definecolor{radarPurple}{HTML}{6B4E9E}
\definecolor{radarOrange}{HTML}{E8734A}
\definecolor{radarBlue}{HTML}{3B7A9E}
\definecolor{radarGreen}{HTML}{4A9B6E}
\definecolor{radarBG}{HTML}{FAFAF8}
\definecolor{radarGrey}{HTML}{5A6472}
\definecolor{kuldefault}{HTML}{00407a}
\definecolor{kuldark}{HTML}{123C75}    % heading background
\definecolor{kullight}{HTML}{F0F7FA}   % block background
\definecolor{kulbg}{RGB}{29,141,176}   % background fill
\definecolor{kulblue}{HTML}{00407a}

\colorlet{blocktitlebgcolor}{radarBlue}
\colorlet{blocktitlefgcolor}{white}
\colorlet{blockbodybgcolor}{white}
\colorlet{backgroundcolor}{radarBG}
\colorlet{framecolor}{radarNavy}

% Semantic aliases used across the library so individual icon
% files never hard-code a raw color -- only these roles.
\colorlet{radarProcessing}{radarBlue}   % radar-domain blocks
\colorlet{deepLearning}{radarGreen}     % NN / learning blocks
\colorlet{resultsColor}{radarOrange}    % results / metrics
\colorlet{geometryColor}{radarPurple}   % geometry / point-cloud blocks
\colorlet{iconStroke}{radarNavy}
\colorlet{iconFaint}{radarGrey!35}
\colorlet{cardBorder}{radarNavy}
\colorlet{cardBG}{white}

% ------------------------------------------------------------
% 2. GLOBAL PGFKEYS  (per-module color role, size, etc.)
% ------------------------------------------------------------
% Every icon/card macro in this library accepts a `color=` key
% that maps to one of the semantic roles above, so a whole
% figure can be recolored by module without touching the icons.
\pgfkeys{
  /radarlib/.is family, /radarlib,
  default/.style = {
    color       = radarProcessing,
    scale       = 1,
    width       = 2.1cm,
  },
  color/.estore in      = \radarlibColor,
  scale/.estore in      = \radarlibScale,
  width/.estore in      = \radarlibWidth,
}
\pgfkeys{/radarlib, default}

% ------------------------------------------------------------
% 3. FONTS / TEXT STYLES
% ------------------------------------------------------------
\tikzset{
  card title/.style   = {font=\sffamily\bfseries\footnotesize, text=blocktitlefgcolor},
  card subtitle/.style = {font=\sffamily\footnotesize, text=radarGrey},
  metric number/.style = {font=\sffamily\bfseries\Large, text=radarNavy},
  metric label/.style  = {font=\sffamily\footnotesize, text=radarGrey},
}
% Plain-text equivalents of the styles above, for use *inside*
% already-open node text (where a tikz style key can't be
% invoked directly) -- e.g. {\metricNumberFont 73.8\%}.
\newcommand{\cardTitleFont}{\sffamily\bfseries\footnotesize\color{blocktitlefgcolor}}
\newcommand{\cardSubtitleFont}{\sffamily\footnotesize\color{radarGrey}}
\newcommand{\metricNumberFont}{\sffamily\bfseries\Large\color{radarNavy}}
\newcommand{\metricLabelFont}{\sffamily\footnotesize\color{radarGrey}}

% ------------------------------------------------------------
% 4. FLOW ARROW  (the connector used between every pipeline stage)
% ------------------------------------------------------------
% Usage:  \draw[flow] (a) -- (b);
%         \draw[flow, color=deepLearning] (a) -- (b);
\tikzset{
  flow/.style = {
    -{Latex[length=2.6mm, width=2mm]},
    line width=0.9pt,
    draw=radarNavy,
    shorten <=1pt,
    shorten >=1pt,
  },
  flow dashed/.style = {
    flow,
    dash pattern=on 2.5pt off 2pt,
  },
}

% Convenience macro so a simple horizontal/vertical chain reads
% declaratively: \PipelineArrow{a}{b}
\newcommand{\PipelineArrow}[3][flow]{\draw[#1] (#2) -- (#3);}

% ------------------------------------------------------------
% 5. PROCESS CARD
% ------------------------------------------------------------
% The workhorse of the whole poster: an icon on top, a colored
% title bar on the bottom, rounded corners, subtle border.
% Built on the `rectangle split` shape so the icon and title
% live in one node (keeps alignment automatic).
%
% Usage inside a tikzpicture:
%   \ProcessCard{a}{Raw Radar}{\RadarCube}
%   \ProcessCard[right=1.4cm of a, titlecolor=deepLearning]{b}{Range FFT}{\RangeFFT}
%
% Args: #1 = extra node options (positioning etc, optional)
%       #2 = node name
%       #3 = title text
%       #4 = icon content (TikZ draw commands, drawn in a small
%            nested picture with origin at the card center)
\tikzset{
  process card/.style = {
    rectangle split,
    rectangle split parts=2,
    rectangle split part align={center, center},
    draw=cardBorder,
    line width=0.8pt,
    rounded corners=4pt,
    text width=2.1cm,
    align=center,
    inner sep=5pt,
    drop shadow={shadow xshift=0.6pt, shadow yshift=-0.6pt, opacity=0.25},
    rectangle split part fill={cardBG, radarProcessing},
  },
  % Safe way to recolor a single card's title bar, e.g.:
  %   \ProcessCard[card color=deepLearning]{b}{Range FFT}{...}
  % (Do not set `fill=` directly on a process card -- it conflicts
  % with `rectangle split part fill` and blanks the node.)
  card color/.style = {rectangle split part fill={cardBG, #1}},
}

\newcommand{\ProcessCard}[4][]{%
  \node[process card, #1] (#2) {%
    \nodepart{one}
      \begin{tikzpicture}[baseline, remember picture, every node/.style={shape=rectangle}]
        #4
      \end{tikzpicture}
    \nodepart{two}
      \textcolor{blocktitlefgcolor}{\sffamily\bfseries\footnotesize #3}
  };%
}

% ------------------------------------------------------------
% 6. METRIC BOX  (the big "73.8% IoU" style callouts)
% ------------------------------------------------------------
% Usage: \MetricBox{a}{0,0}{73.8\%}{IoU}
\newcommand{\MetricBox}[4]{%
  \node[align=center, inner sep=2pt] (#1) at (#2) {%
    {\metricNumberFont #3}\\[1pt]
    {\metricLabelFont #4}%
  };%
}

% ------------------------------------------------------------
% 7. LEGEND BOX  (small colored swatch + label, for figure legends)
% ------------------------------------------------------------
% Usage: \LegendSwatch{0,0}{radarBlue}{Radar processing}
\newcommand{\LegendSwatch}[3]{%
  \begin{scope}[shift={(#1)}]
    \node[fill=#2, rounded corners=1.5pt, minimum width=10pt, minimum height=10pt, anchor=west] (sw) {};
    \node[anchor=west, font=\sffamily\footnotesize, text=radarGrey, right=4pt of sw] {#3};
  \end{scope}
}

% ------------------------------------------------------------
% 8. SECTION HEADER  (poster column headers, e.g. "Contribution 1")
% ------------------------------------------------------------
\newcommand{\SectionHeader}[2][radarProcessing]{%
  \node[fill=#1, text=white, rounded corners=2pt, inner sep=6pt,
        font=\sffamily\bfseries\normalsize, minimum width=4cm] {#2};%
}

% ------------------------------------------------------------
% 9. RESULTS TABLE  (palette-matched booktabs-style table)
% ------------------------------------------------------------
% Not a TikZ node -- a small LaTeX environment wrapper so results
% tables share the same fonts/colors/weight as the figures. Needs
% \usepackage{booktabs} and \usepackage[table]{xcolor} in the
% poster/paper preamble.
%
% Usage:
%   \begin{ResultsTable}{lcc}
%     Method          & Supervised & Few-shot \\
%     TCPCN           & 96.2       & 35.1     \\
%     PCTtrans        & 99.0       & 32.0     \\
%     \BestRow Ours + Pretrain & 94.4 & 70.0  \\
%   \end{ResultsTable}
\newenvironment{ResultsTable}[1]{%
  \renewcommand{\arraystretch}{1.35}%
  \begin{tabular}{#1}
  \toprule
}{%
  \bottomrule
  \end{tabular}%
}
% Header row: white bold text on the radarProcessing band. Formatting
% resets after \midrule so it doesn't bleed into the data rows below.
\newcommand{\HeaderRow}[1]{%
  \rowcolor{radarProcessing}\bfseries\color{white}#1\\
  \midrule
  \normalfont\color{black}%
}
% Highlight a single data row (e.g. your best result) in results color.
\newcommand{\BestRow}{\rowcolor{resultsColor!25}}

% ------------------------------------------------------------
% 10. LEGEND BOX (complete)
% ------------------------------------------------------------
% Usage: \LegendBox{x}{y}{color}{label}
\newcommand{\LegendBox}[4]{%
  \begin{scope}[shift={(#1,#2)}]
    \node[fill=#3, rounded corners=1.5pt, minimum width=8pt, minimum height=8pt, anchor=west] (sw) at (0,0) {};
    \node[anchor=west, font=\sffamily\footnotesize, text=radarGrey, right=6pt of sw] {#4};
  \end{scope}
}